\section{Methodology}
\label{sec:method}

Our goal is to decide, for a skill $S$ that improves a task, whether the improvement reflects
\emph{genuine new knowledge} the base model lacked or \emph{re-elicitation} of latent knowledge.
We first build a substrate where the ground truth of this question is known
(\secref{sec:substrate}), then define the discriminator (\secref{sec:end}), and finally
describe three experiments (\secref{sec:experiments}).

\subsection{A substrate where novelty is known by construction}
\label{sec:substrate}

We use \emph{counterfactual} tasks whose regime we can dial, all verified deterministically
against computable ground truth so no LLM-as-judge is needed. \Tabref{tab:substrate}
summarizes the four settings. The elicitation settings (\arithbten, \vocabes) require a
procedure that is provably in-distribution, so a skill can only re-elicit it. The
novelty-procedure settings (\arithbnine, \arithbeleven) use the canonical counterfactual of
base-$b$ arithmetic~\citep{reasoningreciting}: the correct procedure is verifiable but
out-of-distribution and unreliable. The novelty-information setting (\vocabzeth) is the key
construction: the 24-word invented-language glossary is generated from a fixed random seed at
run time, so the mapping \emph{cannot} be in any training corpus. It is the cleanest lab analog
of a new declarative fact the model has not memorized, yet answers are exactly checkable by
string match.

\begin{table}[t]
    \centering
    \resizebox{\textwidth}{!}{%
    \begin{tabular}{@{}lllp{6.4cm}@{}}
        \toprule
        \textbf{Setting} & \textbf{Task} & \textbf{Regime} & \textbf{What a skill \emph{could} be} \\
        \midrule
        \arithbten & base-10 3-digit multiplication & elicitation & re-elicit a latent, reliable skill \\
        \vocabes & Spanish$\to$English (24-word glossary) & elicitation & re-elicit words the model already knows \\
        \arithbnine, \arithbeleven & base-9 / base-11 3-digit multiplication & novelty-procedure & scaffold a latent-but-unreliable procedure \\
        \vocabzeth & invented-language (``\zeth'')$\to$English & novelty-information & supply information the model provably cannot have (glossary randomly generated at run time) \\
        \bottomrule
    \end{tabular}
    }
    \caption{The controlled substrate. We dial the regime from \emph{elicitation} (procedure is
    in-distribution, a skill can only re-elicit) to \emph{novelty}. \vocabzeth is the cleanest
    novelty case: its 24 word mappings are randomly generated per run, so they cannot appear in
    any training corpus, yet answers are deterministically checkable.}
    \label{tab:substrate}
\end{table}


\subsection{The Elicitation--Novelty Discriminator (\END)}
\label{sec:end}

Fix the base model and a held-out test set. For a skill $S$ with test accuracy $\mathrm{acc}_S$
against a no-skill baseline $\mathrm{acc}_0$, we define the \emph{normalized gain}
\begin{equation}
g(S) = \frac{\mathrm{acc}_S - \mathrm{acc}_0}{1 - \mathrm{acc}_0},
\label{eq:gain}
\end{equation}
which measures the fraction of the available headroom the skill captures. The \END then
comprises three measurable quantities.

\para{Self-generation utility (\SG).} We ask the model to write its own skill from the task
description alone --- no answers --- and reuse it on the test set, giving a self-generated skill
$S_{\text{self}}$. We report $\SG = g(S_{\text{self}}) / g(S_{\text{correct}})$. A high \SG
means the model can produce its own working skill, so the knowledge was already latent
(elicitation); a low \SG means it cannot (candidate novelty).

\para{Length-controlled gain (\LCG).} A skill may help simply by adding context length or
shaping the prompt. We construct a length-matched irrelevant memo $S_{\text{len}}$ and report
$\LCG = g(S_{\text{correct}}) - g(S_{\text{len}})$. Real content yields $\LCG>0$; a pure
prompt-shape effect yields $\LCG\approx0$.

\para{Verifier-only self-evolution.} We run a Voyager/STaR-style loop that refines the model's
own skill from experience, under two feedback signals: \emph{verifier-only} (pass/fail) versus
\emph{with-answer} (ground-truth revealed). If self-evolution can reach the correct-skill
ceiling under verifier-only feedback, the knowledge was latent; if it stalls under
verifier-only feedback but unlocks with ground-truth answers, the knowledge came from outside.

\para{Predicted signature.} \emph{Elicitation} $\Rightarrow$ small $g$, high \SG, $\LCG\approx0$,
verifier-only self-evolution succeeds. \emph{Novelty} $\Rightarrow$ large $g$, low \SG,
$\LCG>0$, verifier-only self-evolution stalls while an external signal unlocks it. The
falsifiable claim is that self-evolution and self-generation sit at the elicitation floor
precisely because \SG is low in the novelty regime: one cannot self-evolve knowledge one does
not have.

\subsection{Three experiments}
\label{sec:experiments}

\para{Exp~1 --- Discriminator.} $5$ settings $\times$ $2$ models $\times$ $4$ conditions
(no-skill, correct-skill, self-generated-skill, length-matched control) $\times$ $N=60$
held-out items. Self-generated skills are written by the model unaided and then reused on the
test set. Models: \claude and \gpt via OpenRouter.

\para{Exp~2 --- Self-evolution loop.} A Voyager/STaR-style loop ($K=5$) that refines the
model's own skill from experience under verifier-only versus with-answer feedback. We track the
learning curve on a disjoint 40-item test set for \arithbnine, \vocabzeth, and \vocabes.
Model: \claude.

\para{Exp~3 --- Internalization probe.} We test the widely proposed ``fine-tune the skill back
into the weights'' idea on a local GPU. We LoRA fine-tune \qwen to put the skill into weights,
then test \emph{without} the skill in context on a held-out split. We compare two fine-tuning
modes: \emph{doc} (train on the skill document as prose) versus \emph{examples} (train on
input$\to$output pairs). Settings: \vocabzeth and \arithbnine.

\subsection{Reproducibility}
\label{sec:repro}

Seeds are fixed at 42 and temperature is 0 (deterministic); all API responses are cached
($\approx$3{,}000 calls). Verifiers are deterministic. We report bootstrap 95\% confidence
intervals (5{,}000 resamples), two-proportion $z$-tests, and Cohen's $h$ effect sizes. LoRA
uses rank $16$, $\alpha=32$, $3$ epochs, learning rate $2\times10^{-4}$, bf16, on a single RTX
A6000; the software stack is Python 3.12, torch 2.5.1+cu124, transformers 4.56.2, and peft.
Our original novelty-information task was a substitution cipher, but provider guardrails on
OpenRouter blocked it as jailbreak-adjacent (\texttt{finish\_reason: content\_filter}); we
replaced it with the mathematically equivalent but innocuous invented-glossary task, and
document this because it affects reproducibility.
